% Auto-generated annex from experiment metadata JSON
\subsection*{Anexo tecnico de metricas por modelo}
\addcontentsline{toc}{subsection}{Anexo tecnico de metricas por modelo}
Este anexo consolida metricas detalladas por modelo y experimento, variables descartadas en pruning, hiperparametros optimos de la fase 5, importancia de caracteristicas en fases 2 y 3, y el arbol de decision final.

\subsubsection*{Experimento 1 (Baseline): metricas detalladas}
\begin{table}[H]
\centering
\caption{Resumen de discriminacion y calibracion - Experimento 1 (Baseline)}
\scriptsize
\begin{tabular}{|l|c|c|c|}
\hline
\textbf{Modelo} & \textbf{AUC-ROC CV} & \textbf{Brier} & \textbf{Threshold} \\ \hline
catboost & 0.6153 & 0.2134 & 0.2200 \\ \hline
decision\_tree & 0.5516 & 0.3877 & 0.2000 \\ \hline
lightgbm & 0.5943 & 0.2394 & 0.2000 \\ \hline
logistic\_regression\_lasso & 0.6055 & 0.2122 & 0.2500 \\ \hline
naive\_bayes & 0.5713 & 0.5993 & 0.2500 \\ \hline
random\_forest & 0.5822 & 0.2454 & 0.2000 \\ \hline
xgboost & 0.5950 & 0.2476 & 0.2000 \\ \hline
\end{tabular}
\end{table}
\begin{table}[H]
\centering
\caption{Resumen de punto operativo - Experimento 1 (Baseline)}
\scriptsize
\begin{tabular}{|l|c|c|c|c|}
\hline
\textbf{Modelo} & \textbf{Recall} & \textbf{Specificity} & \textbf{F1} & \textbf{Kappa} \\ \hline
catboost & 0.8090 & 0.2991 & 0.4911 & 0.0806 \\ \hline
decision\_tree & 0.4111 & 0.6946 & 0.3995 & 0.1042 \\ \hline
lightgbm & 0.6897 & 0.4330 & 0.4771 & 0.0993 \\ \hline
logistic\_regression\_lasso & 0.8196 & 0.2991 & 0.4960 & 0.0883 \\ \hline
naive\_bayes & 0.8594 & 0.1564 & 0.4713 & 0.0111 \\ \hline
random\_forest & 0.7215 & 0.4018 & 0.4827 & 0.0977 \\ \hline
xgboost & 0.7003 & 0.4606 & 0.4925 & 0.1315 \\ \hline
\end{tabular}
\end{table}

\subsubsection*{Experimento 2 (Agregaciones): metricas detalladas}
\begin{table}[H]
\centering
\caption{Resumen de discriminacion y calibracion - Experimento 2 (Agregaciones)}
\scriptsize
\begin{tabular}{|l|c|c|c|}
\hline
\textbf{Modelo} & \textbf{AUC-ROC CV} & \textbf{Brier} & \textbf{Threshold} \\ \hline
catboost & 0.6235 & 0.2121 & 0.2200 \\ \hline
decision\_tree & 0.5523 & 0.3878 & 0.2000 \\ \hline
lightgbm & 0.5969 & 0.2389 & 0.2000 \\ \hline
logistic\_regression\_lasso & 0.6044 & 0.2127 & 0.2500 \\ \hline
naive\_bayes & 0.5804 & 0.5992 & 0.2000 \\ \hline
random\_forest & 0.5912 & 0.2436 & 0.2000 \\ \hline
xgboost & 0.5954 & 0.2466 & 0.2000 \\ \hline
\end{tabular}
\end{table}
\begin{table}[H]
\centering
\caption{Resumen de punto operativo - Experimento 2 (Agregaciones)}
\scriptsize
\begin{tabular}{|l|c|c|c|c|}
\hline
\textbf{Modelo} & \textbf{Recall} & \textbf{Specificity} & \textbf{F1} & \textbf{Kappa} \\ \hline
catboost & 0.8011 & 0.3091 & 0.4907 & 0.0826 \\ \hline
decision\_tree & 0.4111 & 0.6909 & 0.3979 & 0.1002 \\ \hline
lightgbm & 0.6976 & 0.4318 & 0.4808 & 0.1045 \\ \hline
logistic\_regression\_lasso & 0.8064 & 0.2929 & 0.4880 & 0.0738 \\ \hline
naive\_bayes & 0.8462 & 0.1740 & 0.4705 & 0.0142 \\ \hline
random\_forest & 0.6897 & 0.4230 & 0.4736 & 0.0908 \\ \hline
xgboost & 0.6790 & 0.4618 & 0.4817 & 0.1158 \\ \hline
\end{tabular}
\end{table}

\subsubsection*{Experimento 3 (Categorizaciones): metricas detalladas}
\begin{table}[H]
\centering
\caption{Resumen de discriminacion y calibracion - Experimento 3 (Categorizaciones)}
\scriptsize
\begin{tabular}{|l|c|c|c|}
\hline
\textbf{Modelo} & \textbf{AUC-ROC CV} & \textbf{Brier} & \textbf{Threshold} \\ \hline
catboost & 0.6134 & 0.2148 & 0.2200 \\ \hline
decision\_tree & 0.5550 & 0.3876 & 0.2000 \\ \hline
lightgbm & 0.5934 & 0.2392 & 0.2000 \\ \hline
logistic\_regression\_lasso & 0.6037 & 0.2125 & 0.2500 \\ \hline
naive\_bayes & 0.5741 & 0.5972 & 0.2900 \\ \hline
random\_forest & 0.5794 & 0.2453 & 0.2000 \\ \hline
xgboost & 0.5950 & 0.2476 & 0.2000 \\ \hline
\end{tabular}
\end{table}
\begin{table}[H]
\centering
\caption{Resumen de punto operativo - Experimento 3 (Categorizaciones)}
\scriptsize
\begin{tabular}{|l|c|c|c|c|}
\hline
\textbf{Modelo} & \textbf{Recall} & \textbf{Specificity} & \textbf{F1} & \textbf{Kappa} \\ \hline
catboost & 0.8143 & 0.2979 & 0.4932 & 0.0835 \\ \hline
decision\_tree & 0.4164 & 0.6934 & 0.4031 & 0.1079 \\ \hline
lightgbm & 0.7003 & 0.4293 & 0.4813 & 0.1045 \\ \hline
logistic\_regression\_lasso & 0.8090 & 0.3004 & 0.4915 & 0.0816 \\ \hline
naive\_bayes & 0.8568 & 0.1640 & 0.4722 & 0.0145 \\ \hline
random\_forest & 0.7029 & 0.4005 & 0.4728 & 0.0823 \\ \hline
xgboost & 0.7003 & 0.4606 & 0.4925 & 0.1315 \\ \hline
\end{tabular}
\end{table}

\subsubsection*{Experimento 4 (Pruning): metricas detalladas}
\begin{table}[H]
\centering
\caption{Resumen de discriminacion y calibracion - Experimento 4 (Pruning)}
\scriptsize
\begin{tabular}{|l|c|c|c|}
\hline
\textbf{Modelo} & \textbf{AUC-ROC CV} & \textbf{Brier} & \textbf{Threshold} \\ \hline
catboost & 0.6209 & 0.2129 & 0.2300 \\ \hline
decision\_tree & 0.5549 & 0.3844 & 0.2000 \\ \hline
lightgbm & 0.5969 & 0.2389 & 0.2000 \\ \hline
logistic\_regression\_lasso & 0.6043 & 0.2116 & 0.2400 \\ \hline
naive\_bayes & 0.5829 & 0.2681 & 0.2000 \\ \hline
random\_forest & 0.5974 & 0.2419 & 0.2000 \\ \hline
xgboost & 0.6040 & 0.2423 & 0.2000 \\ \hline
\end{tabular}
\end{table}
\begin{table}[H]
\centering
\caption{Resumen de punto operativo - Experimento 4 (Pruning)}
\scriptsize
\begin{tabular}{|l|c|c|c|c|}
\hline
\textbf{Modelo} & \textbf{Recall} & \textbf{Specificity} & \textbf{F1} & \textbf{Kappa} \\ \hline
catboost & 0.8064 & 0.3292 & 0.4996 & 0.1023 \\ \hline
decision\_tree & 0.3899 & 0.6909 & 0.3813 & 0.0798 \\ \hline
lightgbm & 0.6976 & 0.4318 & 0.4808 & 0.1045 \\ \hline
logistic\_regression\_lasso & 0.8515 & 0.2628 & 0.4988 & 0.0832 \\ \hline
naive\_bayes & 0.4350 & 0.6596 & 0.4034 & 0.0908 \\ \hline
random\_forest & 0.7162 & 0.4205 & 0.4865 & 0.1094 \\ \hline
xgboost & 0.6950 & 0.4631 & 0.4906 & 0.1295 \\ \hline
\end{tabular}
\end{table}

\subsubsection*{Experimento 5 (Hiperparametros): metricas detalladas}
\begin{table}[H]
\centering
\caption{Resumen de discriminacion y calibracion - Experimento 5 (Hiperparametros)}
\scriptsize
\begin{tabular}{|l|c|c|c|}
\hline
\textbf{Modelo} & \textbf{AUC-ROC CV} & \textbf{Brier} & \textbf{Threshold} \\ \hline
catboost & 0.6202 & 0.2151 & 0.3600 \\ \hline
decision\_tree & 0.6227 & 0.2142 & 0.2200 \\ \hline
lightgbm & 0.6338 & 0.2143 & 0.3200 \\ \hline
logistic\_regression\_lasso & 0.6133 & 0.2129 & 0.3000 \\ \hline
naive\_bayes & 0.5840 & 0.2627 & 0.2000 \\ \hline
random\_forest & 0.6356 & 0.2314 & 0.4200 \\ \hline
xgboost & 0.6391 & 0.2672 & 0.4800 \\ \hline
\end{tabular}
\end{table}
\begin{table}[H]
\centering
\caption{Resumen de punto operativo - Experimento 5 (Hiperparametros)}
\scriptsize
\begin{tabular}{|l|c|c|c|c|}
\hline
\textbf{Modelo} & \textbf{Recall} & \textbf{Specificity} & \textbf{F1} & \textbf{Kappa} \\ \hline
catboost & 0.8329 & 0.3104 & 0.5056 & 0.1068 \\ \hline
decision\_tree & 0.8329 & 0.3304 & 0.5122 & 0.1227 \\ \hline
lightgbm & 0.8992 & 0.2616 & 0.5191 & 0.1159 \\ \hline
logistic\_regression\_lasso & 0.8382 & 0.2804 & 0.4984 & 0.0872 \\ \hline
naive\_bayes & 0.4748 & 0.6308 & 0.4207 & 0.0989 \\ \hline
random\_forest & 0.8249 & 0.3630 & 0.5196 & 0.1434 \\ \hline
xgboost & 0.8647 & 0.2854 & 0.5118 & 0.1100 \\ \hline
\end{tabular}
\end{table}

\subsubsection*{Fase 4 (Pruning): variables descartadas}
Se descartaron 7 variables con criterio de ocurrencia minima (10).
\begin{table}[H]
\centering
\caption{Variables descartadas en feature pruning}
\begin{tabular}{|l|}
\hline
\textbf{Variable} \\ \hline
ENFERMEDAD\_CORONARIA \\ \hline
FALLA\_CARDIACA \\ \hline
INMUNOSUPRESION \\ \hline
HIPERTENSION\_PULMONAR \\ \hline
EPOC \\ \hline
ASMA \\ \hline
ENF\_CEREBROVASCULAR \\ \hline
\end{tabular}
\end{table}

\subsubsection*{Fase 5 (Hiperparametros): parametros optimos por modelo}
\begin{table}[H]
\centering
\caption{Parametros optimos (optimized\_params) en Experimento 5}
\scriptsize
\begin{tabular}{|l|p{10.5cm}|}
\hline
\textbf{Modelo} & \textbf{Parametros optimos} \\ \hline
catboost & verbose=0, iterations=100, learning\_rate=0.01, depth=5, l2\_leaf\_reg=1, auto\_class\_weights=None \\ \hline
decision\_tree & max\_depth=7, min\_samples\_split=10, min\_samples\_leaf=20, criterion=gini, class\_weight=None \\ \hline
lightgbm & verbosity=-1, n\_estimators=100, learning\_rate=0.01, max\_depth=5, num\_leaves=31, min\_child\_samples=20, subsample=1, colsample\_bytree=1, scale\_pos\_weight=2 \\ \hline
logistic\_regression\_lasso & penalty=l2, solver=lbfgs, max\_iter=10000, C=0.001, class\_weight=None \\ \hline
naive\_bayes & var\_smoothing=0.1 \\ \hline
random\_forest & n\_jobs=-1, n\_estimators=300, max\_depth=5, min\_samples\_split=10, min\_samples\_leaf=20, max\_features=None, class\_weight=balanced \\ \hline
xgboost & verbosity=0, n\_estimators=500, max\_depth=7, learning\_rate=0.01, subsample=1, colsample\_bytree=0.6, scale\_pos\_weight=3, gamma=5, reg\_alpha=0, reg\_lambda=1 \\ \hline
\end{tabular}
\end{table}

\subsubsection*{Importancia de caracteristicas por modelo (Fases 2 y 3)}
Para cada modelo se reporta el top 3 de importancia cuando el artefacto esta disponible en \texttt{plot\_metadata.json}. En modelos sin importancia embebida se marca como No aplica.
\begin{table}[H]
\centering
\caption{Top 3 de importancia de caracteristicas en Experimentos 2 y 3}
\scriptsize
\begin{tabular}{|l|l|p{8.5cm}|}
\hline
\textbf{Experimento} & \textbf{Modelo} & \textbf{Top 3 caracteristicas (pct)} \\ \hline
Experimento 2 (Agregaciones) & catboost & HB\_PREQX (26.83\%), EDAD2 (21.34\%), EDAD (20.75\%) \\ \hline
Experimento 2 (Agregaciones) & decision\_tree & HB\_PREQX (35.30\%), EDAD (22.49\%), EDAD2 (16.24\%) \\ \hline
Experimento 2 (Agregaciones) & lightgbm & HB\_PREQX (44.60\%), EDAD (39.07\%), GENERO (3.43\%) \\ \hline
Experimento 2 (Agregaciones) & logistic\_regression\_lasso & No aplica (modelo sin importancia embebida en el artefacto) \\ \hline
Experimento 2 (Agregaciones) & naive\_bayes & No aplica (modelo sin importancia embebida en el artefacto) \\ \hline
Experimento 2 (Agregaciones) & random\_forest & HB\_PREQX (40.91\%), EDAD (20.64\%), EDAD2 (19.43\%) \\ \hline
Experimento 2 (Agregaciones) & xgboost & SANGRADO\_MAYOR (26.68\%), HIPERTENSION (7.11\%), ACT\_FISICA\_METS (6.08\%) \\ \hline
Experimento 3 (Categorizaciones) & catboost & EDAD (36.92\%), HB\_PREQX (24.89\%), SANGRADO\_MAYOR (9.03\%) \\ \hline
Experimento 3 (Categorizaciones) & decision\_tree & EDAD (38.69\%), HB\_PREQX (36.01\%), GENERO (5.87\%) \\ \hline
Experimento 3 (Categorizaciones) & lightgbm & HB\_PREQX (45.03\%), EDAD (42.43\%), GENERO (3.37\%) \\ \hline
Experimento 3 (Categorizaciones) & logistic\_regression\_lasso & No aplica (modelo sin importancia embebida en el artefacto) \\ \hline
Experimento 3 (Categorizaciones) & naive\_bayes & No aplica (modelo sin importancia embebida en el artefacto) \\ \hline
Experimento 3 (Categorizaciones) & random\_forest & EDAD (38.72\%), HB\_PREQX (36.36\%), SANGRADO\_MAYOR (4.20\%) \\ \hline
Experimento 3 (Categorizaciones) & xgboost & SANGRADO\_MAYOR (28.88\%), CANCER\_ACTIVO (9.05\%), HIPOTIROIDISMO (7.64\%) \\ \hline
\end{tabular}
\end{table}

\subsubsection*{Arbol de decision final (Experimento 5)}
Se incluye el arbol completo generado en la carpeta de plots del modelo decision\_tree del Experimento 5.
\begin{figure}[H]
\centering
\includegraphics[width=0.98\textwidth]{../experiments/exp_05_hyperparameter_search/models/decision_tree/plots/decision_tree_full.png}
\caption{Arbol de decision completo - Experimento 5 (decision\_tree)}
\label{fig:anexo-exp5-decision-tree-full}
\end{figure}
